% ═══════════════════════════════════════════════════════════════════
% §6  공기열원 히트펌프 급탕기 (동적 모델)
% Source: AirSourceHeatPumpBoiler.py
% ═══════════════════════════════════════════════════════════════════
\section{공기열원 히트펌프 급탕기}
\label{sec:ashpb}

본 절은 \texttt{AirSourceHeatPumpBoiler.py}에 구현된 동적 ASHPB 모델을 기술한다. 증기압축 냉매 사이클, 물리 기반 열교환기, Differential Evolution 최적화기, 히스테리시스 제어 저탕조를 통합한다.

\subsection{가정}
\begin{enumerate}
  \item 냉매 사이클은 각 시간 단계 내에서 준정상상태로 운전된다.
  \item 냉매 물성은 CoolProp(실제유체 상태방정식)으로 계산한다.
  \item 과열도 $\Delta T_\text{sh}$와 과냉도 $\Delta T_\text{sc}$는 일정하다.
  \item 압축기는 배제 체적 기반 질량류량 모델과 등엔트로피 효율을 따른다.
  \item 응축기는 침수 코일($UA\cdot\Delta T$ 모델, LMTD 아님)이다.
  \item 증발기는 공기측 에너지 수지(교차 유동)를 사용한다.
\end{enumerate}

\subsection{냉매 사이클 상태점}

포화 온도는 열원/열원 조건에 의해 결정된다:
\begin{align}
  T_\text{eva,sat} &= T_0 - \Delta T_\text{ref,evap}  \label{eq:ashp-Teva}\\
  T_\text{cond,sat} &= T_\text{w,tank} + \Delta T_\text{ref,cond} \label{eq:ashp-Tcond}
\end{align}
여기서 $\Delta T_\text{ref,evap}$, $\Delta T_\text{ref,cond}$는 Differential Evolution 최적화 변수이다.

압축기 입구 (과열):
\begin{equation}\label{eq:ashp-cmp-in}
  T_\text{cmp,in} = T_\text{eva,sat} + \Delta T_\text{sh}
  \,,\quad
  h_\text{cmp,in} = h(T_\text{cmp,in},\,P_\text{eva})
\end{equation}

압축기 출구 (등엔트로피 + 효율):
\begin{equation}\label{eq:ashp-cmp-out}
  h_\text{cmp,out} = h_\text{cmp,in} + \frac{h_\text{cmp,out,s} - h_\text{cmp,in}}{\eta_\text{isen}}
\end{equation}

\subsection{열교환기 모델}

\paragraph{응축기 (침수 코일).}
\begin{equation}\label{eq:ashp-Qcond}
  Q_\text{cond} = UA_\text{cond}\,\Delta T_\text{ref,cond}
\end{equation}

\paragraph{증발기 (공기측 에너지 수지).}
\begin{equation}\label{eq:ashp-Qevap}
  Q_\text{evap} = \rho_a\,c_a\,\dot{V}_\text{fan}\,(T_\text{a,in} - T_\text{a,out})
\end{equation}
팬 유량 $\dot{V}_\text{fan}$은 공기측 용량과 냉매측 증발 부하 $Q_\text{ref,evap} = \dot{m}_\text{ref}(h_\text{eva,out} - h_\text{exp,out})$를 일치시켜 결정한다.

\paragraph{증발기 열관류율 스케일링.}
증발기의 총합 열관류율($UA$)은 공기 풍량($\dot{V}_\text{fan}$)의 변동에 따라 동적으로 스케일링된다. 본 모델은 계산 효율성을 위해 열교환기 형상을 단순화하였으며, 공기측 열전달 계수($h$)와 풍량($\dot{V}_\text{fan}$)의 관계를 다음의 척도법(Scaling law)을 적용하여 모델링하였다:
\begin{equation}\label{eq:ashp-UA}
  UA = UA_\text{ref} \left( \frac{\dot{V}_\text{fan}}{\dot{V}_\text{ref}} \right)^n
\end{equation}
여기서 지수 $n=0.65$는 Wang et al. (2000)~\cite{wang2000heat}이 제시한 평면 핀-관(Plain fin-and-tube) 열교환기 상관식에 근거하여 다음과 같이 유도된다.

첫째, Colburn $j$-팩터의 정의에 따라 열전달 계수 $h$는 최대 풍속 $w_\text{max}$($\propto \dot{V}_\text{fan}$)와 $j$의 곱에 비례한다 ($h \propto j \cdot \dot{V}_\text{fan}$). 둘째, Wang의 상관식에 따르면 일반적인 공조용 열교환기 운전 범위($1,000 < Re < 10,000$, 4열 이상 관)에서 평면 핀의 $j$-팩터는 레이놀즈 수($Re$)의 지수 함수로 근사되며, $j \propto Re^{-0.35} \propto \dot{V}_\text{fan}^{-0.35}$의 비례 관계를 따른다. 이를 $h$에 대한 비례식에 대입하면 다음과 같은 지수 법칙(Power-law)이 성립한다:
\begin{equation}\label{eq:ashp-h-scaling}
  h \propto \dot{V}_\text{fan} \cdot \dot{V}_\text{fan}^{-0.35} = \dot{V}_\text{fan}^{0.65}
\end{equation}
이러한 $\dot{V}_\text{fan}^{0.65}$ 스케일링은 핀 사이의 좁은 간격으로 인해 층류와 난류가 혼재된 전이 영역(Transition region)의 핀-관 열전달 특성을 반영한 것으로, 일반적인 완전 발달 난류 관내 유동 상관식(예: Dittus-Boelter 식, $n=0.8$)보다 낮은 지수를 가지게 된다.

\subsection{압축기 모델}

배제 체적으로부터의 질량류량:
\begin{equation}\label{eq:ashp-mdot}
  \dot{m}_\text{ref} = \rho_\text{cmp,in}\,V_\text{disp}\,\frac{n_\text{cmp}}{60}
\end{equation}
여기서 $n_\text{cmp}$ [RPM]은 응축기 열수지에 의해 결정된다.

소비 전력:
\begin{equation}\label{eq:ashp-Ecmp}
  E_\text{cmp} = \dot{m}_\text{ref}\,(h_\text{cmp,out} - h_\text{cmp,in})
\end{equation}

\subsection{Differential Evolution 최적화}

\paragraph{목적함수.}
총 소비 전력 최소화:
\begin{equation}\label{eq:ashp-obj}
  \min_{\Delta T_\text{ref,cond},\,\Delta T_\text{ref,evap}} \;
  E_\text{tot} = E_\text{cmp} + E_\text{fan,ou}
\end{equation}

\paragraph{제약 조건.}
\begin{itemize}
  \item 열수지: $Q_\text{cond} = \dot{m}_\text{ref}(h_\text{cmp,out} - h_\text{exp,in})$
  \item 팬 유량 범위: $0.1\,\dot{V}_\text{design} \le \dot{V}_\text{fan} \le \dot{V}_\text{design}$
  \item 압축기 회전수 범위: $n_\text{min} \le n_\text{cmp} \le n_\text{max}$
\end{itemize}

\subsection{저탕조 에너지 밸런스}
\label{sec:ashpb-tank}

저탕조는 가변 질량의 집중정전용량(lumped capacitance) 모델로 취급한다.
지배 상미분방정식은 다음과 같다:
\begin{equation}\label{eq:ashp-tank-ode}
  C \frac{\dd T}{\dd t} = \dot{Q}_\text{gain} - UA_\text{tank}\,(T - T_0)
\end{equation}
여기서 $C = c_w \rho_w V_\text{tank} \cdot \ell$은 유효 열용량(저탕조 수위 $\ell \in (0,1]$에 의해 스케일링),
$\dot{Q}_\text{gain}$은 열원 항(응축기, UV 램프, STC, 급수 엔탈피)의 합이다.

\paragraph{완전 음해법(Fully Implicit) 이산화.}
모든 온도 의존 항을 $n{+}1$ 시점에서 평가하며, 질량--에너지 연립 시스템을
Newton--Raphson(\texttt{scipy.optimize.fsolve})으로 동시에 풀어 구한다.

\textbf{에너지 잔차:}
\begin{equation}\label{eq:ashp-implicit-energy}
  r_E = C_\text{full}\,\ell^{n+1}\,T^{n+1}
      - C_\text{full}\,\ell^{n}\,T^{n}
      - \Delta t\,\bigl[
          \dot{Q}_\text{HP} + E_\text{UV} + E_\text{pump}
          + \dot{Q}_\text{STC}(T^{n+1})
          + \dot{Q}_\text{refill}(T^{n+1})
          - UA\,(T^{n+1} - T_0)
        \bigr]
\end{equation}
여기서
$\dot{Q}_\text{refill}(T) = c_w \rho_w \dot{V}_\text{in}(T_\text{refill} - T)$,
$\dot{Q}_\text{STC}(T)$는 탱크 온도 $T$에서 평가된 순 STC 이득이다.

\textbf{질량 잔차:}
\begin{equation}\label{eq:ashp-implicit-mass}
  r_M = \ell^{n+1} - \ell^{n}
      - \frac{\dot{V}_\text{in} - \dot{V}_\text{out}(T^{n+1})}{V_\text{full}}\,\Delta t
\end{equation}
여기서 $\dot{V}_\text{out}(T) = \alpha(T)\,\dot{V}_\text{mix}$이며,
믹싱 비율은
$\alpha(T) = \text{clip}\bigl(\frac{T_\text{mix}-T_\text{in}}{T - T_\text{in}},\;0,\;1\bigr)$이다.

미지수 $[T^{n+1},\,\ell^{n+1}]$는 $\mathbf{r} = \mathbf{0}$을
\texttt{fsolve}로 풀어 결정한다.
이 방법은 무조건 안정하며, 가변 탱크 질량과 온도의 강한 연성을 처리할 수 있다.

\subsection{엑서지 분석}

서브시스템 엑서지 수지는 입력--파괴--출력 패턴을 따른다~\cite{tsatsaronis2007}:

\paragraph{냉매 루프.}
\begin{equation}\label{eq:ashp-Xr}
  X_\text{cmp} - X_\text{c,r} = X_\text{r,iu} + X_\text{r,ou}
\end{equation}

\paragraph{실내기 (응축기측).}
\begin{equation}\label{eq:ashp-Xiu}
  X_\text{r,iu} - X_\text{c,iu} = X_\text{w,out} - X_\text{w,in}
\end{equation}

\paragraph{전체 엑서지 효율.}
\begin{equation}\label{eq:ashp-etaX}
  \eta_{X,\text{sys}} = \frac{X_\text{w,out} - X_\text{w,in}}{X_\text{cmp} + X_\text{fan,iu} + X_\text{fan,ou}}
\end{equation}
